% ======================================================================
% Chapter 18: Operator Products and Effective Vertices
% (book pages 599-647; PDF pages 622-673)
% ======================================================================
\chapter{算符乘积与有效顶角}

我们在第~17 章中对 QCD 的分析建立在渐近自由原理之上，它告诉我们，具有大动量转移的强相互作用过程可以可靠地用弱耦合微扰论来处理。然而，到目前为止，我们对 QCD 中的重整化群这一更强有力的工具还很少使用。在本章中，我们将推出 Callan-Symanzik 方程在 QCD 中的一些推论。我们将看到，渐近自由理论具有它们自己特有的标度行为，其修正具有动量标度对数的反常幂次的形式。尽管这些修正一般弱于第~13 章所研究的标量场理论中的修正，它们仍然对强相互作用有重要的定性影响。

我们从考虑 QCD 中质量项的标度律开始，直接采用我们在第~12.4 和~12.5 节中用来描述 $\phi^4$ 理论质量项的形式体系。然而，其他应用则需要更强有力的理论工具，即\emph{算符乘积展开}。第~18.3 节引入量子场论中算符乘积的一般描述，并解释这样的算符乘积如何受到 Callan-Symanzik 方程的约束。最后两节利用这一工具来发展对深度非弹性散射和 QCD 中其他硬过程的新观点。

\section{夸克质量参数的重整化}

到目前为止，我们一直假定夸克质量足够小，以至于在高能过程中可以忽略它们。即使对轻夸克 $u$、$d$、$s$，这也不总是一个恰当的假定；对较重的夸克 $c$、$b$、$t$，质量可能具有非常重要的影响。然而，由于孤立的夸克并不存在，不可能无歧义地定义夸克的质量。在下面的讨论中，我们将把夸克质量视为 QCD 微扰论的一个参数，由某个重整化标度 $M$ 处的重整化方案来定义。

因为我们像定义耦合常数那样、通过重整化约定来定义夸克质量，所以我们应当预期这个参数会按照重整化群演化而跑动，从而不同的质量参数值适用于不同的过程。我们说，我们最初的方案导致一个\emph{有效夸克质量}，它取决于对其进行计算的动量标度。在本节中，我们将推出这一有效质量对动量标度的领头依赖关系。

有效质量项的基本形式体系已在第~12.5 节中给出。要把质量项加到 QCD 拉格朗日量中，我们必须首先通过在标度 $M$ 处的重整化方案来定义质量算符 $(\bar{q}q)$。然后我们可以通过在拉格朗日量中加入如下项来定义夸克质量
\begin{equation}
  \Delta\mathcal{L}_m = -m\,(\bar{q}q)_M .
  \label{eq:18.1}
\end{equation}
在这个讨论中，我们将假定夸克质量 $m$ 足够小，以至于我们只需保留 $m$ 的领头阶项。我们还为简单起见假定，我们只对一种夸克味有这样的质量项。

在零质量极限下，算符 $(\bar{q}q)$ 与夸克场的 Green 函数
\begin{equation}
  G^{(n;k)}(x_1, \ldots, x_n; y_1, \ldots, y_n; z_1, \ldots, z_k)
  = \bigl\langle q(x_1)\cdots q(x_n)\bar{q}(y_1)\cdots\bar{q}(y_n)\,
    \bar{q}q(z_1)\cdots\bar{q}q(z_k) \bigr\rangle ,
  \label{eq:18.2}
\end{equation}
满足 Callan-Symanzik 方程
\begin{equation}
  \Bigl[ M\frac{\partial}{\partial M} + \beta\frac{\partial}{\partial g}
    + 2n\gamma + k\gamma_{\bar{q}q} \Bigr]\ G^{(n;k)}\bigl( \{x_i\}; g, M
    \bigr) = 0 ,
  \label{eq:18.3}
\end{equation}
其中 $\gamma$ 是夸克场的反常维数，$\gamma_{\bar{q}q}$ 是算符 $\bar{q}q$ 的反常维数。如果我们按照 \eqref{eq:18.1} 把质量项包含在拉格朗日量中，则 $n$ 个夸克场和 $n$ 个反夸克场的 Green 函数满足
\begin{equation}
  \Bigl[ M\frac{\partial}{\partial M} + \beta\frac{\partial}{\partial g}
    + 2n\gamma + m\frac{\partial}{\partial m} \Bigr]\ G^{(n)}\bigl(
      \{x_i\}, \{y_i\}; g, m, M \bigr) = 0 .
  \label{eq:18.4}
\end{equation}
对 $m$ 的导数对质量算符被使用的次数计数。在第~12.5 节中，我们把具有质量量纲的变量 $m$ 换成一个无量纲变量。然而，在 QCD 中，把有量纲的参数 $m$ 视为耦合常数同样方便。于是 Callan-Symanzik 方程的解将包含一个跑动质量参数 $\bar{m}(Q)$，它取决于 Green 函数的一个典型动量 $Q$。这个参数被定义为与式~\eqref{eq:12.126} 类似的一个重整化群方程的解。对于这种情况，方程为
\begin{equation}
  \frac{d}{d\log(Q/M)}\ \bar{m} = \gamma_{\bar{q}q}(\bar{g})\cdot\bar{m} ,
  \label{eq:18.5}
\end{equation}
初始条件为
\begin{equation}
  \bar{m}(M) = m .
  \label{eq:18.6}
\end{equation}
量 $\bar{m}(Q)$ 就是\emph{有效质量}，应当用它来计算动量转移为 $Q$ 的夸克产生或散射过程的质量效应。

为了显式地计算 $\bar{m}(Q)$，我们需要算出质量算符的反常维数 $\gamma_{\bar{q}q}$。这可以按照第~12.4 节所解释的那样来做。我们通过如下方案显式地定义算符的归一化：重整化夸克场之间的 $(\bar{q}q)$ 顶角函数应当满足
\begin{equation}
  \frac{}{p^2-q^2=(p+q)^2=-M^2} = 1 .
  \label{eq:18.7}
\end{equation}
为了保持 \eqref{eq:18.7}，我们将需要一个具有算符插入结构的抵消项顶角 $\delta_{\bar{q}q}$。于是，与式~\eqref{eq:12.112} 类似，反常维数在单圈阶由下式给出
\begin{equation}
  \gamma_{\bar{q}q} = \frac{\partial}{\partial\log M}\Bigl(
    -\delta_{\bar{q}q} + \delta_2 \Bigr) ,
  \label{eq:18.8}
\end{equation}
其中 $\delta_2$ 是夸克场强重整化的抵消项，定义于图~\ref{fig:16.8}。规范不变算符 $(\bar{q}q)$ 的关联函数是规范不变的，因此该函数的 Callan-Symanzik 方程中的各项之和必须是规范不变的结果。由于 $\beta(g)$ 的领头系数与规范及其他约定无关，由 \eqref{eq:18.3} 可知 $\gamma_{\bar{q}q}$ 的领头系数也与约定无关。抵消项 $\delta_2$ 和 $\delta_{\bar{q}q}$ 都依赖于规范。这一论证表明，规范依赖必须在 \eqref{eq:18.8} 中消去。在下面的计算以及本章中的其他反常维数计算中，我们将自始至终在 Feynman-'t Hooft 规范中计算。

我们已于第~16.4 节在 Feynman-'t Hooft 规范中计算了抵消项 $\delta_2$ 的发散部分。对 QCD 求结果 \eqref{eq:16.77} 中的群论因子，我们得到
\begin{equation}
  \delta_2 = \frac{4}{3}\ \frac{g^2}{(4\pi)^2}\ \frac{\Gamma(2-d/2)}
    {(M^2)^{2-d/2}} .
  \label{eq:18.9}
\end{equation}
为了计算 $\delta_{\bar{q}q}$，我们必须算出顶角 \eqref{eq:18.7} 的单圈修正。它由如下图形给出
\begin{equation}
  \text{(one-loop correction to } \bar{q}q \text{ vertex)}
  = -ig^2\ \int \frac{d^dk}{(2\pi)^d}\ \frac{1}{(k+q)^2}\ \frac{1}{k^2}\
    \frac{1}{(k-p)^2}\ \cdots ,
  \label{eq:18.10}
\end{equation}
在这个图形的表达式中，因子 1 表示 $\bar{q}q$ 算符插入。在相应的夸克数流 $j^\mu = \bar{q}\gamma^\mu q$ 重整化的图形中，这个因子将被 $\gamma^\mu$ 取代。由于我们只需要顶角重整化 \eqref{eq:18.10} 的发散部分，我们可以把被积函数近似为它在大的 $k$ 处的值。于是这个图形变为
\begin{equation}
  \text{(diagram)} \sim ig^2\ \frac{2}{(4\pi)^2}\ \Gamma(2-d/2)\ \cdots .
  \label{eq:18.11}
\end{equation}
为了保持归一化条件 \eqref{eq:18.7}，我们必须加入抵消项
\begin{equation}
  \delta_{\bar{q}q} = \frac{8}{3}\ \frac{g^2}{(4\pi)^2}\ \frac{\Gamma(2-d/2)}
    {(M^2)^{2-d/2}} .
  \label{eq:18.12}
\end{equation}
综合 \eqref{eq:18.8}、\eqref{eq:18.9} 和 \eqref{eq:18.12}，我们得到
\begin{equation}
  \gamma_{\bar{q}q} = \frac{4}{3}\ \frac{g^2}{(4\pi)^2} .
  \label{eq:18.13}
\end{equation}

正如我们在上一段中所指出的，夸克数流的反常维数可以通过非常相似的计算得到。这将对我们的形式体系提供一个良好的检验，因为正如我们在式~\eqref{eq:12.110} 上方所论证的，守恒流由它的积分、即守恒荷，无歧义地归一化，因此其反常维数必为零。如果我们把 \eqref{eq:18.10} 中的 1 换成 $j^\mu$，并使用同一套近似来化简积分，我们在分子中得到 Dirac 矩阵结构
\begin{equation}
  \gamma^\rho\gamma^\mu\gamma^\nu\gamma_\rho\gamma^\sigma\gamma_\mu
  \to \frac{1}{2}(-2)^2\,\gamma^\nu\,\cdots ,
  \label{eq:18.14}
\end{equation}
于是，代替 \eqref{eq:18.12}，我们需要抵消项
\begin{equation}
  \delta_j = \frac{4}{3}\ \frac{g^2}{(4\pi)^2}\ \frac{\Gamma(2-d/2)}
    {(M^2)^{2-d/2}} .
  \label{eq:18.15}
\end{equation}
把这一结果与 \eqref{eq:18.9} 结合，我们得到
\begin{equation}
  \gamma_j = 0 ,
  \label{eq:18.16}
\end{equation}
与我们的普遍论证相符。

如果我们把 \eqref{eq:18.11} 中的 gamma 函数换成 $\log(\Lambda^2/Q^2)$ 的显式因子，然后用抵消项 \eqref{eq:18.12} 减去发散，我们发现顶角图形的行为为
\begin{equation}
  \text{(vertex)} \sim \frac{4}{3}\ \frac{g^2}{8\pi^2}\
    \log\frac{M^2}{Q^2} .
  \label{eq:18.17}
\end{equation}
这个图形在小的外动量处给出增强。这种增强的一部分与外夸克场的（与规范有关的）重标度相联系；关系式 \eqref{eq:18.8} 告诉我们如何提取这个对数中与有效质量的规范不变增强相联系的部分。于是，到 $\alpha_s$ 阶，
\begin{equation}
  \bar{m}(Q) = m\ \Bigl[ 1 + \frac{4}{3}\frac{\alpha_s}{\pi}
    \log\frac{M}{Q} + \cdots \Bigr] .
  \label{eq:18.18}
\end{equation}

\begin{figure}[t]
  \centering
  \includegraphics[width=0.55\textwidth]{fig18.1.png}
  \caption{给出夸克有效质量动量依赖的领头对数贡献的图形。}
  \label{fig:18.1}
\end{figure}

为了更精确地计算有效质量的动量依赖，我们必须把计算中的另外两个特征考虑进去。首先，量 $\alpha_s\log(M^2/Q^2)$ 可能变得量级为 1，在这种情况下，我们必须把所有形如 $(\alpha_s\log(M^2/Q^2))^n$ 的领头对数项都考虑进去。这类贡献来自图~\ref{fig:18.1} 中所示的所有图形。其次，耦合常数 $\alpha_s$ 本身是动量标度的函数，这使来自小 $Q$ 的贡献得到进一步的增强。这两种效应都可以通过求解重整化群方程 \eqref{eq:18.5} 而正确地计及。到 $g^2$ 的领头阶，这个方程具有显式形式
\begin{equation}
  \frac{d}{d\log(Q/M)}\ \bar{m} = \frac{4}{3}\ \frac{\bar{g}^2}{16\pi^2}\,
    \bar{m} .
  \label{eq:18.19}
\end{equation}
把 $g$ 的重整化群方程的解以形式 \eqref{eq:17.17} 代入，我们得到
\begin{equation}
  \bar{m}(Q) = \bar{m}(\mu)\ \Bigl[ \frac{\log(\mu^2/\Lambda^2)}
    {\log(Q^2/\Lambda^2)} \Bigr]^{4/(3b_0)} .
  \label{eq:18.20}
\end{equation}
有效夸克质量随动量标度的增大而缓慢减小。这是 QCD 的一个具体的、可检验的预言：夸克质量的跑动。

\section{弱相互作用的 QCD 重整化}

在本节中，我们将把重整化群形式体系应用于弱相互作用。夸克的弱相互作用涉及四费米子算符，即我们将在第~20 章中详细讨论的那种类型。我们将看到，QCD 对这些算符的修正可能很重要，并且它们会导致算符混合。

我们回忆，$u$、$d$ 和 $s$ 夸克的弱相互作用由有效拉格朗日量描述
\begin{equation}
  \Delta\mathcal{L} = -\frac{G_F}{\sqrt{2}}\ \bigl( \bar{u}_L\gamma^\mu d_L
    \bigr)\ \bigl( \bar{e}_L\gamma_\mu\nu_{eL} \bigr)\ \cos\theta_c
    + \cdots ,
  \label{eq:18.28}
\end{equation}
其中 $\theta_c$ 是 Cabibbo 角。这一项允许 $s$ 夸克通过过程 $s\to uud$ 衰变。类似地，\eqref{eq:18.28} 的转动产生有效相互作用
\begin{equation}
  \Delta\mathcal{L} = -\frac{4G_F}{\sqrt{2}}\ \sin\theta_c\,
    \bigl( \bar{s}_L\gamma^\mu u_L \bigr)\ \bigl( \bar{u}_L\gamma_\mu d_L
      \bigr) ,
  \label{eq:18.32}
\end{equation}
它导致衰变 $s\to u\bar{u}d$。这些弱相互作用过程分别被称为非轻子衰变过程和半轻子衰变过程。类似的表达式适用于其他重夸克。

鉴于式~(18.31) 和 \eqref{eq:18.32} 在基本层次上描述 $s$ 夸克的弱相互作用耦合，我们现在讨论这些耦合受 QCD 对数的修正。我们已在前一节看到，QCD 修正对增强底层拉格朗日量中夸克质量项的强度有深刻的影响。我们现在要研究弱相互作用的强度是否也能受到类似的增强。

\begin{figure}[t]
  \centering
  \includegraphics[width=0.5\textwidth]{fig18.2.png}
  \caption{半轻子弱相互作用顶角强度的 QCD 修正。}
  \label{fig:18.2}
\end{figure}

我们首先考虑半轻子弱相互作用算符 \eqref{eq:18.32}。轻子费米子双线性不受 QCD 影响，因此该算符的 QCD 增强与其夸克分量
\begin{equation}
  \bar{u}_L\gamma^\mu s_L .
  \label{eq:18.33}
\end{equation}
的增强完全相同。然而，这个算符是一个流，因此 $\gamma = 0$。就图形而言，图~\ref{fig:18.2} 中所示图形产生的对数增强被夸克场强重整化抵消，正如我们在第~18.1 节对流的顶角的讨论中已经看到的那样。左手投影算符 $\tfrac{1}{2}(1-\gamma^5)$ 与图形对易，对最终结果没有影响。同样的评述适用于联结 $u$ 和 $d$ 夸克的半轻子弱相互作用。对那种情况，它意味着式~\eqref{eq:17.35} 中给出的深度非弹性中微子散射截面的归一化不受 QCD 对数的影响。

\begin{figure}[t]
  \centering
  \includegraphics[width=0.55\textwidth]{fig18.3.png}
  \caption{非轻子弱相互作用顶角强度的 QCD 修正。}
  \label{fig:18.3}
\end{figure}

然而，在非轻子弱相互作用的情形下，QCD 的效应并不那么简单。让我们首先计算给出弱相互作用顶角式~(18.31) 重整化领头修正的费曼图，然后在稍后的阶段建立这些结果的重整化群解释。

在 $\alpha_s$ 阶，非轻子弱相互作用顶角从图~\ref{fig:18.3} 中所示图形得到修正。注意第一个图形正是半轻子情形中所发现的流重整化。第二个图形给出第二个夸克流的类似重整化。在 $\gamma$ 的计算中，这两个贡献抵消了四个夸克场场强重整化的贡献。图~\ref{fig:18.3} 中其余四个图形是新的贡献，它们给出可能很大的重标度因子。

我们现在计算这些图形，从图~\ref{fig:18.3} 中的第三个图形开始。与第~18.1 节的计算一样，我们感兴趣的是与远大于外动量的圈动量 $k$ 值相联系的对数发散贡献。提取这一贡献的最简单方法是在零外动量近似下计算每个图形。在写下这些图形的表达式时，我们将省略前置因子
\begin{equation}
  \frac{4G_F}{\sqrt{2}}\ \cos\theta_c\sin\theta_c .
  \label{eq:18.34}
\end{equation}
我们将保留夸克场来表示外态，从而使我们的最终表达式具有重标度算符的形式。

使用这种记号，图~\ref{fig:18.3} 中第三个图形的值为
\begin{equation}
  \text{(third diagram)} = ig^2\ \int \frac{d^dk}{(2\pi)^d}\ \frac{1}{k^4}\
    \bigl( \bar{d}_L\gamma^\rho\gamma^\lambda\gamma^\sigma t^a u_L \bigr)\
    \bigl( \bar{u}_L\gamma^\rho\gamma_\lambda\gamma_\sigma t^a s_L \bigr) .
  \label{eq:18.35}
\end{equation}
利用 $k$ 积分的对称性，我们提取发散部分：
\begin{equation}
  \begin{aligned}
    \text{(third diagram)} &= ig^2\ \int \frac{d^dk}{(2\pi)^d}\ \frac{1}{k^4}\
      \bigl( \bar{d}_L\gamma^\rho\gamma^\lambda\gamma^\sigma t^a u_L \bigr)\
      \bigl( \bar{u}_L\gamma^\rho\gamma_\lambda\gamma_\sigma t^a s_L \bigr)\\
    &= -g^2\ \frac{\Gamma(2-d/2)}{(4\pi)^2}\ \bigl( \bar{d}_L\gamma^\lambda
      \gamma^\mu\gamma_\lambda t^a u_L \bigr)\ \bigl( \bar{u}_L\gamma_\mu
        t^a s_L \bigr)\ \cdots .
  \end{aligned}
  \label{eq:18.36}
\end{equation}
为了把夸克场的乘积放到更熟悉的形式中，我们应用第~3.4 节末尾讨论的 Fierz 变换。如果不存在色矩阵 $t^a$，费米子场的乘积就恰好是 \eqref{eq:3.82} 中出现的那个，我们将得到
\begin{equation}
  \bigl( \bar{d}_L\gamma^\mu\gamma^\lambda\gamma_\mu u_L \bigr)\
    \bigl( \bar{u}_L\gamma_\lambda s_L \bigr) = 16\ \bar{d}_L\gamma^\mu u_L\
    \bar{u}_L\gamma_\mu s_L .
  \label{eq:18.37}
\end{equation}
矩阵 $t^a$ 重新引导夸克场的色量子数。为了澄清这一点，我们需要恒等式 \eqref{eq:3.77} 的色类比。为找到这一恒等式，考虑色不变量
\begin{equation}
  (t^a)_{ij}(t^a)_{k\ell} .
  \label{eq:18.38}
\end{equation}
指标 $i$、$k$ 按照色的 3 表示变换；指标 $j$、$\ell$ 按照 $\bar{3}$ 变换。因此，\eqref{eq:18.38} 必须是收缩这些指标的两种可能方式的线性组合，
\begin{equation}
  A\,\delta_{i\ell}\delta_{kj} + B\,\delta_{ij}\delta_{k\ell} .
  \label{eq:18.39}
\end{equation}
常数 $A$ 和 $B$ 可以通过把 \eqref{eq:18.38} 和 \eqref{eq:18.39} 与 $\delta_{ij}$ 及与 $\delta_{jk}$ 收缩，并调整 $A$ 和 $B$ 使 \eqref{eq:18.39} 的收缩满足恒等式
\begin{equation}
  \mathrm{tr}[t^a](t^a)_{k\ell} = 0;\qquad
  (t^at^a)_{i\ell} = \frac{4}{3}\delta_{i\ell} .
  \label{eq:18.40}
\end{equation}
来确定。这给出恒等式
\begin{equation}
  (t^a)_{ij}(t^a)_{k\ell} = \frac{1}{2}\Bigl( \delta_{i\ell}\delta_{kj}
    - \frac{1}{3}\delta_{ij}\delta_{k\ell} \Bigr) .
  \label{eq:18.41}
\end{equation}
对基表示中 $SU(N)$ 的生成元，也有类似的关系成立，在那种情形下 $(1/3)$ 被 $(1/N)$ 取代。

把 \eqref{eq:18.41} 插入 \eqref{eq:18.36}，我们发现该恒等式的第一项产生一个新的四费米子算符，
\begin{equation}
  \bigl( \bar{d}_{Li}\gamma^\mu u_{Lj} \bigr)\ \bigl( \bar{u}_{Lj}
    \gamma_\mu s_{Li} \bigr) ,
  \label{eq:18.42}
\end{equation}
其中 $i$、$j$ 是色指标。应用 \eqref{eq:18.37} 中的 Fierz 重排，然后再应用附加的重排 \eqref{eq:3.79}，我们可以把这个算符化为如下形式
\begin{equation}
  16\ \bigl( \bar{d}_{Li}\gamma^\mu u_{Lj} \bigr)\ \bigl( \bar{u}_{Lj}
    \gamma_\mu s_{Li} \bigr) = 16\ \bigl( \bar{u}_{Lj}\gamma^\mu u_{Lj}
    \bigr)\ \bigl( \bar{d}_{Li}\gamma_\mu s_{Li} \bigr) .
  \label{eq:18.43}
\end{equation}
\eqref{eq:3.79} 中的负号被交换费米子场顺序所产生的一个负号补偿。最终结果是色单态夸克流的乘积；然而，这些流中的场与原始算符中的场结合方式不同。

我们对这个图形计算的最终结果是
\begin{equation}
  \text{(third diagram)} = -4g^2\ \frac{\Gamma(2-d/2)}{(4\pi)^2}\ \Bigl[
    \bar{d}_L\gamma^\mu u_L\ \bar{u}_L\gamma_\mu s_L
    - \frac{1}{3}\bar{d}_L\gamma^\mu u_L\ \bar{u}_L\gamma_\mu s_L
    \Bigr] .
  \label{eq:18.44}
\end{equation}
图~\ref{fig:18.3} 的第四个图形给出完全相同的贡献。

图~\ref{fig:18.3} 中最后两个图形的计算非常类似。第五个图形给出
\begin{equation}
  \text{(fifth diagram)} = ig^2\ \int \frac{d^dk}{(2\pi)^d}\ \frac{1}{k^4}\
    \bigl( \bar{d}_L\gamma^\rho\gamma^\lambda\gamma^\sigma t^a u_L \bigr)\
    \bigl( \bar{u}_L\gamma_\rho\gamma_\lambda\gamma_\sigma t^a s_L \bigr)\
    \cdots = +g^2\ \frac{\Gamma(2-d/2)}{(4\pi)^2}\ \bigl( \bar{d}_L\gamma^\mu
    t^a u_L \bigr)\ \bigl( \bar{u}_L\gamma_\mu t^a s_L \bigr) .
  \label{eq:18.45}
\end{equation}
利用 Fierz 恒等式 \eqref{eq:3.79}，四费米子算符可以如下化简：
\begin{equation}
  \begin{aligned}
    \bigl( \bar{d}_L\gamma^\mu t^a u_L \bigr)\ \bigl( \bar{u}_L\gamma_\mu
      t^a s_L \bigr) &= +\bigl( \bar{s}_L\gamma^\mu\gamma^\lambda\gamma_\mu
      t^a s_L \bigr)\ \bigl( \bar{u}_L\gamma_\lambda t^a u_L \bigr)\\
    &= +(-2)^2\ \bigl( \bar{d}_L\gamma^\mu t^a u_L \bigr)\ \bigl(
      \bar{u}_L\gamma_\mu t^a s_L \bigr) .
  \end{aligned}
  \label{eq:18.46}
\end{equation}
同样，我们必须用恒等式 \eqref{eq:18.41} 化简色矩阵的乘积，并且该恒等式的第一项再次需要一次额外的 Fierz 变换。最终结果是
\begin{equation}
  \text{(fifth diagram)} = +g^2\ \frac{\Gamma(2-d/2)}{(4\pi)^2}\ \Bigl[
    \bar{d}_L\gamma^\mu u_L\ \bar{u}_L\gamma_\mu s_L
    - \frac{1}{3}\bar{d}_L\gamma^\mu u_L\ \bar{u}_L\gamma_\mu s_L
    \Bigr] .
  \label{eq:18.47}
\end{equation}
图~\ref{fig:18.3} 中最后一个图形给出相同的贡献。这四个图形贡献之和为
\begin{equation}
  -3g^2\ \frac{\Gamma(2-d/2)}{(4\pi)^2}\ \Bigl[ \bar{d}_L\gamma^\mu u_L\
    \bar{u}_L\gamma_\mu s_L - \frac{1}{3}\bar{d}_L\gamma^\mu u_L\
      \bar{u}_L\gamma_\mu s_L \Bigr] .
  \label{eq:18.48}
\end{equation}

提取图~\ref{fig:18.3} 中诸图形的紫外发散部分是计算弱相互作用顶角 Callan-Symanzik $\gamma$ 函数的正式方案的一部分。然而，在此处稍作停顿，问一问这个发散的物理意义是有益的。如果我们用含有 $W$ 玻色子的底层理论来计算，图~\ref{fig:18.3} 中的图形不会发散。在把弱相互作用写成一个有效的局域顶角时，我们把 $W$ 玻色子传播子近似为一个常数，假定它所携带的动量 $k$ 远小于 $m_W$：
\begin{equation}
  \frac{1}{k^2-m_W^2} \to -\frac{1}{m_W^2}\ \cdots .
  \label{eq:18.49}
\end{equation}
我们用来计算有效顶角 QCD 修正的近似只在 $k^2 \ll m_W^2$ 的积分区域中有效。在这个区域之外，我们必须使用完整的 $W$ 传播子；这将在分母中引入一个额外的 $k^2$ 因子并使积分收敛。因此，在 QCD 修正的直接计算中，图~\ref{fig:18.3} 计算中的紫外发散项将被在 $m_W$ 处截断的对数取代。对数的下限由外动量确定。在 $K$ 介子——含有 $s$ 夸克的最轻强子——的衰变中，这些动量量级为 $m_K$。因此，\eqref{eq:18.48} 中给出的修正应当通过如下替换来计算
\begin{equation}
  \frac{\Gamma(2-d/2)}{(4\pi)^2} \to \frac{1}{4\pi}\ \log\frac{m_W^2}{m_K^2}
  .
  \label{eq:18.50}
\end{equation}
按照这种解释，我们可以把 \eqref{eq:18.48} 改写为对领头阶弱相互作用顶角的 $\alpha_s$ 阶修正。这一修正的效应是弱相互作用算符的重标度和修改：
\begin{equation}
  \begin{aligned}
    \bar{d}_L\gamma^\mu u_L\ \bar{u}_L\gamma_\mu s_L
    &\to \Bigl( 1 + \frac{\alpha_s}{4\pi}\log\frac{m_W^2}{m_K^2} \Bigr)\
      \bar{d}_L\gamma^\mu u_L\ \bar{u}_L\gamma_\mu s_L\\
    &\quad - 3\ \frac{\alpha_s}{4\pi}\log\frac{m_W^2}{m_K^2}\
      \bar{u}_L\gamma^\mu u_L\ \bar{d}_L\gamma_\mu s_L .
  \end{aligned}
  \label{eq:18.51}
\end{equation}
注意 QCD 修正不仅重标度了原始算符的归一化，还引入了一个具有不同结构的新算符。这个计算使第~12.4 节中引入的思想具体化，即改变局域算符归一化的图形也可能把具有相同维数和量子数的不同算符混合在一起。

由于 \eqref{eq:18.50} 中对数的值约为 10，领头 QCD 修正的大小量级为 1，因此高阶修正是重要的。为了对领头对数修正求和，我们回到重整化群分析。为清晰起见，定义
\begin{equation}
  \mathcal{O}_1 = \bar{d}_L\gamma^\mu u_L\ \bar{u}_L\gamma_\mu s_L;
  \qquad \mathcal{O}_2 = \bar{u}_L\gamma^\mu u_L\ \bar{d}_L\gamma_\mu s_L .
  \label{eq:18.52}
\end{equation}
这些算符的反常维数矩阵可以从单圈计算中读出。求解重整化群方程，我们发现非轻子衰变的有效哈密顿量包含这两个算符与跑动耦合之比幂次的组合，导致 $\Delta I = 1/2$ 振幅的显著增强。这一效应称为 $\Delta I = 1/2$ 规则的增强，是 QCD 的经典成就之一。

\section{算符乘积展开}

在前一节中，我们遇到了局域算符乘积的一个重要例子：弱相互作用的四费米子算符。在本节中，我们将发展一个处理量子场论中算符乘积的一般形式体系，即由 Wilson 提出的\emph{算符乘积展开}（OPE）。这个形式体系将使我们能够系统地分析算符乘积的矩阵元在短距离极限下的行为。

考虑两个算符 $\mathcal{O}_1(x)$ 和 $\mathcal{O}_2(0)$ 的乘积与一组场 $\phi(y_i)$ 的普通 Green 函数：
\begin{equation}
  G_{12}\bigl( x; y_1, \ldots, y_m \bigr)
  = \bigl\langle \mathcal{O}_1(x)\mathcal{O}_2(0)\ \phi(y_1)\cdots\phi(y_m)
    \bigr\rangle .
  \label{eq:18.63}
\end{equation}
对小的 $x$，这个 Green 函数是奇异的，因为这两个算符几乎在同一点。奇异性来自两个算符内部场之间的收缩。然而，奇异的收缩只能涉及一组有限的量子数：所得局域算符必须具有乘积 $\mathcal{O}_1\mathcal{O}_2$ 的整体对称性量子数，但除此之外不受限制。把这个算符写成来自标准基的算符的线性组合是有用的。这个线性组合中的系数可以依赖于分离 $x$。通常，量子场论中算符的乘积是奇异的，因此很可能某些系数会随着 $x\to 0$ 而具有奇异性。综合这些观察，Wilson 提出，算符乘积的效应可以通过把 \eqref{eq:18.63} 中的算符乘积替换为局域算符的线性组合来计算，
\begin{equation}
  \mathcal{O}_1(x)\mathcal{O}_2(0) \to \sum_n C_{12n}(x)\,
    \mathcal{O}_n(0) ,
  \label{eq:18.64}
\end{equation}
其中系数 $C_{12n}(x)$ 是 c 数函数。这个算符乘积展开（OPE）将只依赖于算符 $\mathcal{O}_1$、$\mathcal{O}_2$ 及它们的分离，并且与 Green 函数中出现的其他场的同一性和位置无关。

展开 \eqref{eq:18.64} 意味着 Green 函数 \eqref{eq:18.63} 可以按如下方式对小的 $x$ 展开：
\begin{equation}
  G_{12}\bigl( x; y_1, \ldots, y_m \bigr) = \sum_n C_{12n}(x)\,
    G_n\bigl( y_1, \ldots, y_m \bigr) ,
  \label{eq:18.65}
\end{equation}
其中
\begin{equation}
  G_n\bigl( y_1, \ldots, y_m \bigr) = \bigl\langle \mathcal{O}_n(0)\,
    \phi(y_1)\cdots\phi(y_m) \bigr\rangle ,
  \label{eq:18.66}
\end{equation}
而且对 $x$ 的所有依赖现在都由 OPE 系数函数承载。在前一节的例子中，最终振幅以相当复杂的方式依赖于两个算符的小分离，这是通过式~(18.61) 中系数对 $m_W$ 的依赖实现的。从算符乘积展开的观点来看，这种依赖由系数函数承载，并且当这些系数被计算出来时，对所有矩阵元都被确定了。

在第~18.1 和~18.2 节中，我们用重整化群计算了算符矩阵元的增强或压低因子。因此，自然可以预期，算符乘积系数的形式也由重整化群决定。我们现在将推出这一关系。首先，我们把展开 \eqref{eq:18.64} 更精确地改写。出现在这个关系中的算符必须在某个重整化标度 $M$ 处定义。于是算符乘积展开为：
\begin{equation}
  \bigl[ \mathcal{O}_1(x) \bigr]_M \bigl[ \mathcal{O}_2(0) \bigr]_M
  = \sum_n C_{12n}(x; M)\ \bigl[ \mathcal{O}_n(0) \bigr]_M .
  \label{eq:18.67}
\end{equation}
注意系数函数可以依赖于 $M$，因为它们必须吸收与 $M$ 有关的算符重标度。如果我们用 \eqref{eq:18.67} 的左边来计算 \eqref{eq:18.63}，这个函数满足 Callan-Symanzik 方程
\begin{equation}
  \Bigl[ M\frac{\partial}{\partial M} + \beta\frac{\partial}{\partial g}
    + m\gamma + \gamma_1 + \gamma_2 \Bigr]\ G_{12}\bigl( x; y_1, \ldots,
      y_m; M \bigr) = 0 .
  \label{eq:18.68}
\end{equation}
类似地，当算符 $\mathcal{O}_n$ 在 $M$ 处归一化时，Green 函数 \eqref{eq:18.66} 满足
\begin{equation}
  \Bigl[ M\frac{\partial}{\partial M} + \beta\frac{\partial}{\partial g}
    + m\gamma + \gamma_n \Bigr]\ G_n\bigl( y_1, \ldots, y_m; M \bigr) = 0 .
  \label{eq:18.69}
\end{equation}
把 \eqref{eq:18.68} 应用到 \eqref{eq:18.65} 的右边，我们看到只有当 OPE 系数函数满足 Callan-Symanzik 方程时，这些关系才自洽
\begin{equation}
  \Bigl[ M\frac{\partial}{\partial M} + \beta\frac{\partial}{\partial g}
    + \gamma_1 + \gamma_2 - \gamma_n \Bigr]\ C_{12n}\bigl( x; M \bigr) = 0 .
  \label{eq:18.70}
\end{equation}
我们现在用标准方法求解这个方程。首先，让我们应用量纲分析。如果算符 $\mathcal{O}_1$、$\mathcal{O}_2$、$\mathcal{O}_n$ 具有维数 $d_1$、$d_2$、$d_n$，则系数函数 $C_{12n}(x)$ 必须具有 $(\mathrm{mass})^{d_1+d_2-d_n}$ 的量纲。因此，
\begin{equation}
  C_{12n}(x) = \Bigl( \frac{1}{|x|} \Bigr)^{d_1+d_2-d_n}\
    \bar{C}(xM) ,
  \label{eq:18.71}
\end{equation}
其中 $\bar{C}(xM)$ 是一个无量纲函数。这个函数按照第~12.3 节的方法由 \eqref{eq:18.70} 确定。于是，
\begin{equation}
  C_{12n}(x) = \Bigl( \frac{1}{|x|} \Bigr)^{d_1+d_2-d_n}\ \bar{c}\bigl(
    \bar{g}(1/|x|) \bigr)\ \exp\Bigl[ \int_{1/|x|}^{M} d\log M'\ \bigl(
      \gamma_n - \gamma_1 - \gamma_2 \bigr) \Bigr] ,
  \label{eq:18.72}
\end{equation}
其中 $\bar{c}(g)$ 是分离标度 $1/|x|$ 处跑动耦合常数的无量纲函数。

在重整化群的定点处，$\gamma$ 函数将取确定的值 $\gamma_j^{*} = \gamma_j(g^{*})$。于是，解 \eqref{eq:18.72} 可以计算为
\begin{equation}
  C_{12n}(x) = \Bigl( \frac{1}{|x|} \Bigr)^{d_1+d_2-d_n}\ \bar{c}(g^{*})\
    \exp\Bigl[ \log(xM)\bigl( \gamma_n^{*} - \gamma_1^{*} - \gamma_2^{*}
      \bigr) \Bigr] .
  \label{eq:18.73}
\end{equation}
因此，在这种情形下，
\begin{equation}
  C_{12n}(x) = \bar{c}(g^{*})\ \Bigl( \frac{1}{|x|} \Bigr)^{d_1^{*}+d_2^{*}
    - d_n^{*}} ,
  \label{eq:18.74}
\end{equation}
其中
\begin{equation}
  d_j^{*} = d_j + \gamma_j(g^{*})
  \label{eq:18.75}
\end{equation}
是算符 $\mathcal{O}_j$ 在定点处的真实标度维数。

对渐近自由理论的情形，标度关系以我们在第~18.1 节中推出的那种方式变得复杂。在微扰论的领头阶，三个 $\gamma$ 函数取式~(18.24) 的形式。于是 \eqref{eq:18.72} 的解取如下形式
\begin{equation}
  C_{12n}(x) \sim \Bigl( \frac{1}{|x|} \Bigr)^{d_1+d_2-d_n}\ \Bigl(
    \frac{\log(1/|x|^2\Lambda^2)}{\log(M^2/\Lambda^2)} \Bigr)^{
      (\alpha_n-\alpha_1-\alpha_2)/2b_0} .
  \label{eq:18.76}
\end{equation}
在第~18.2 节的例子中，原始算符是维数为 3 且 $\gamma = 0$ 的流，处于分离 $m_W^{-1}$，最终局域算符的维数为 6。因此，\eqref{eq:18.76} 确实正确地再现了式~(18.61) 对 $m_W$ 的依赖。注意对流的乘积，即归一化与标度无关的流，重整化群依赖不那么复杂。这个特殊情形在算符乘积展开的应用中经常出现。

我们写式~\eqref{eq:18.70} 时没有考虑算符混合。然而，正如我们已经看到的，算符混合往往是 OPE 应用的一个基本组成部分。通过用矩阵值 $\gamma$ 函数改写导致 \eqref{eq:18.70} 的分析，可以直截了当地包含这一效应。例如，在存在算符混合时，$G_n$ 的 Callan-Symanzik 方程将被修改为
\begin{equation}
  \Bigl[ M\frac{\partial}{\partial M} + \beta\frac{\partial}{\partial g}
    + m\gamma \Bigr]\ G_n + \sum_p \gamma_{np}\,G_p\bigl( y_1, \ldots, y_m; M
    \bigr) = 0 .
  \label{eq:18.77}
\end{equation}
做了这些改动后，\eqref{eq:18.70} 变为
\begin{equation}
  \Bigl[ M\frac{\partial}{\partial M} + \beta\frac{\partial}{\partial g}
    \Bigr]\ C_{12n}(x; M) + \gamma_{1k}\,C_{12k}(x; M)
    + \gamma_{2k}\,C_{12n}(x; M) - C_{12k}(x; M)\,\gamma_{kn} = 0 .
  \label{eq:18.78}
\end{equation}
注意前两个 $\gamma$ 矩阵从左边作用在 OPE 系数上，而第三个从右边作用。在流的乘积的情形下，前两个 $\gamma$ 矩阵消失，\eqref{eq:18.78} 简化为
\begin{equation}
  \Bigl[ M\frac{\partial}{\partial M} + \beta\frac{\partial}{\partial g}
    \Bigr]\ C_{12n}(x; M) - C_{12k}(x; M)\,\gamma_{kn} = 0 .
  \label{eq:18.79}
\end{equation}
这个方程将在第~18.5 节的分析中起重要作用。

\section{$e^{+}e^{-}$ 湮灭的算符分析}

不难想象，流彼此置于短距离处的矩阵元与流传递硬动量转移的矩阵元之间存在联系。因此我们可能预期，算符乘积展开的思想将给我们一个新的观点来理解 QCD 中硬散射过程的理论。在本节和下一节中，我们将推出算符乘积展开与第~17 章微扰 QCD 分析的关系。

我们从讨论 $e^{+}e^{-}$ 湮灭为强子的总截面开始。在式~\eqref{eq:17.9} 下方，我们论证了这个总截面可以在 QCD 微扰论中计算，使用与总质心能量标度相对应的 $\alpha_s$ 值。然而，这个论证纯粹是直观的，包含许多逻辑跳跃。在本节中，我们将对同一结论给出一个更严格的论证。

\begin{figure}[t]
  \centering
  \includegraphics[width=0.5\textwidth]{fig18.4.png}
  \caption{其虚部给出 $e^{+}e^{-}\to$ 强子总截面的图形。}
  \label{fig:18.4}
\end{figure}

为了援引算符乘积展开，我们必须把 $e^{+}e^{-}$ 湮灭为强子的总截面写成一个流的乘积的矩阵元。为此，我们使用光学定理把 $e^{+}e^{-}$ 散射的总截面与 $e^{+}e^{-}\to e^{+}e^{-}$ 的前向散射振幅联系起来。忽略电子的质量，我们从式~\eqref{eq:7.49} 看到
\begin{equation}
  \sigma\bigl( e^{+}e^{-} \bigr) = \frac{1}{2s}\ \mathrm{Im}\,
    \mathcal{M}\bigl( e^{+}e^{-}\to e^{+}e^{-} \bigr) .
  \label{eq:18.80}
\end{equation}
为了计算 $e^{+}e^{-}\to$ 强子的截面，我们在计算虚部时只考虑来自强子中间态的贡献。到 $\alpha$ 的领头阶，但对强相互作用到所有阶，这些贡献来自只考虑图~\ref{fig:18.4} 形式的图形，并取真空极化中强子贡献的虚部。

图~\ref{fig:18.4} 中所示图形的值为
\begin{equation}
  i\mathcal{M} = (-ie)^2\ \bar{u}(k)\gamma^\mu v(k_+)\ \frac{1}{q^2}\
    \bigl( i\Pi^{\mu\nu}_H(q) \bigr)\ \frac{1}{q^2}\ \bar{v}(k_+)\gamma^\nu
    u(k) ,
  \label{eq:18.81}
\end{equation}
其中 $s = q^2$，$\Pi^{\mu\nu}_H(q)$ 是真空极化的强子部分。由 Ward 恒等式，这可以写成如下形式
\begin{equation}
  \Pi^{\mu\nu}_H(q) = \bigl( g^{\mu\nu}q^2 - q^\mu q^\nu \bigr)\,
    \Pi_H(q^2) .
  \label{eq:18.82}
\end{equation}
然后通过取虚部得到总强子截面：
\begin{equation}
  \sigma\bigl( e^{+}e^{-}\to\text{hadrons} \bigr) = \frac{4\pi\alpha^2}{s}\
    \mathrm{Im}\,\Pi_H(s) .
  \label{eq:18.83}
\end{equation}
与 \eqref{eq:17.4} 比较，我们看到函数 $\mathrm{Im}\,\Pi_H(s)$ 就是包含 QCD 动力学的量。

真空极化 $\Pi_H(q^2)$ 是电磁流乘积的矩阵元，
\begin{equation}
  \Pi_H(q^2) \propto \int d^4x\ e^{iq\cdot x}\ \bigl\langle J^\mu(x)
    J^\nu(0) \bigr\rangle ,
\end{equation}
其中这些流是夸克的电磁流。对大的 $q^2$，积分由小的 $x$ 主导，我们可以把算符乘积展开应用于流的乘积。两个流的 OPE 包含单位算符、质量算符和算符 $(F^a_{\mu\nu})^2$ 等：
\begin{equation}
  J^\mu(x)J^\nu(0) \sim C_1(x)\,\mathbb{1} + C_{\bar{q}q}(x)\,
    m\,\bar{q}q(0) + C_{G^2}(x)\, (F^a_{\mu\nu})^2(0) + \cdots .
  \label{eq:18.91}
\end{equation}
取真空期望值，领头项给出
\begin{equation}
  \Pi_H(q^2) \approx \Pi_H^{(0)}(q^2) + \mathcal{O}\bigl( 1/q^4 \bigr) ,
  \label{eq:18.94}
\end{equation}
其中 $\Pi_H^{(0)}$ 是领头微扰贡献，在 QCD 微扰论中用标度 $q^2$ 处求值的跑动耦合计算：
\begin{equation}
  \Pi_H^{(0)}(q^2) = \frac{1}{12\pi^2}\ \sum_f Q_f^2\ \Bigl( 1 +
    \frac{\alpha_s(q^2)}{\pi} + \cdots \Bigr) .
  \label{eq:18.95}
\end{equation}
这给出 $e^{+}e^{-}$ 湮灭的总截面，与我们在第~17 章中导出的结果相同。

\begin{figure}[t]
  \centering
  \includegraphics[width=0.5\textwidth]{fig18.6.png}
  \caption{$\Pi_H(q^2)$ 在复 $q^2$ 平面中的解析奇点。}
  \label{fig:18.6}
\end{figure}

\begin{figure}[t]
  \centering
  \includegraphics[width=0.55\textwidth]{fig18.7.png}
  \caption{推导 $\sigma(e^{+}e^{-}\to$ 强子) 的 ITEP 求和规则时涉及的积分围道。}
  \label{fig:18.7}
\end{figure}

算符乘积展开通过 \eqref{eq:18.91} 中更高维数的算符给出超出领头阶结果的额外信息。这些修正可以通过使用色散关系与可观测的截面联系起来。考虑积分
\begin{equation}
  \frac{1}{2\pi i}\oint dq^2\ \frac{1}{(q^2+Q_0^2)^{n+1}}\ \Pi_H(q^2) ,
  \label{eq:18.97}
\end{equation}
对 $n > 1$，在环绕 $q^2 = -Q_0^2$ 的围道上求值。如果我们把围道收缩到极点，我们得到
\begin{equation}
  \frac{1}{2\pi i}\oint dq^2\ \frac{1}{(q^2+Q_0^2)^{n+1}}\ \Pi_H(q^2)
  = \frac{1}{n!}\ \frac{d^n}{dq^{2n}}\ \Pi_H(q^2)\Big|_{q^2=-Q_0^2} ,
\end{equation}
这可以通过从算符乘积关系 \eqref{eq:18.91} 求 $\Pi_H$ 来计算，
\begin{equation}
  \Pi_H(q^2) = -e^2\ \Bigl[ c_1(q^2) + c_{\bar{q}q}(q^2)\,
    \langle 0|m\bar{q}q|0\rangle + c_{G^2}(q^2)\,
      \langle 0|(F^a_{\mu\nu})^2|0\rangle + \cdots \Bigr] .
  \label{eq:18.98}
\end{equation}
另一方面，我们可以通过把围道变形到图~\ref{fig:18.7} 的形式来求该积分的值。由于随着 $q^2\to\infty$，没有任何系数函数比 $(q^2)^0$ 乘对数增长得更快，对 $n > 1$ 可以忽略无穷远处的围道。环绕支割线的围道部分给出
\begin{equation}
  \begin{aligned}
    &= -\frac{1}{4\pi}\ \int_0^{\infty} \frac{dq^2}{(q^2+Q_0^2)^{n+1}}\
      \mathrm{Disc}\,\Pi_H(q^2)\\
    &= \frac{1}{2\pi}\ \int_0^{\infty} \frac{ds}{(s+Q_0^2)^{n+1}}\
      \sigma(s) ,
  \end{aligned}
  \label{eq:18.99}
\end{equation}
其中 $\sigma(s)$ 是 $e^{+}e^{-}\to$ 强子的总截面。令 \eqref{eq:18.97} 和 \eqref{eq:18.99} 相等，我们得到 OPE 系数（在 QCD 微扰论中求值）与可观测截面之间的一系列积分关系。这些关系最初由 Novikov、Shifman、Voloshin、Vainshtein 和 Zakharov 构造，称为\emph{ITEP 求和规则}。%
\footnote{这些求和规则的理论综述见 V.~A.~Novikov、L.~B.~Okun、M.~A.~Shifman、A.~I.~Vainshtein、M.~B.~Voloshin 和 V.~I.~Zakharov, \emph{Phys.\ Repts.} \textbf{41}, 1 (1978)。}

只用 $c_2(q^2)$ 的领头 QCD 表达式来求求和规则，我们得到
\begin{equation}
  \int_0^{\infty} \frac{ds}{(s+Q_0^2)^{n+1}}\ \sigma(s)
  = \frac{2\pi\alpha^2}{3}\ \sum_f Q_f^2\ \frac{1}{Q_0^{2n}}\ \Bigl[ 1 +
    \mathcal{O}\bigl( \alpha_s(Q_0^2) \bigr) + \mathcal{O}\Bigl(
      \frac{\langle G^2\rangle}{Q_0^4} \Bigr) + \cdots \Bigr] .
  \label{eq:18.100}
\end{equation}
领头阶关系与式~\eqref{eq:18.95} 中给出的最低阶截面一致。修正来自 QCD 微扰论的更高阶（其中 $\alpha_s$ 取标度 $Q_0^2$ 处的值）以及 OPE 中更高的算符项。

如果 \eqref{eq:18.100} 中的修正项对 $n$ 一致地收敛到零，我们就可以反演求和规则并从中导出我们的结果 \eqref{eq:18.94}。然而，真实的情形更为微妙。由于 \eqref{eq:18.97} 中的导数强调 $q^2$ 变化更强的项，ITEP 求和规则中的修正项随 $n$ 增大而越来越重要。因此，截面偏离 QCD 微扰论预言的最重要偏差是对这一预言的振荡，这些振荡在低 $n$ 的求和规则中被平均掉。理论与实验的比较见图~\ref{fig:18.8}。在大的 $s$ 处，\eqref{eq:18.94} 相当精确。然而，随着 $s$ 变小，振荡的幅度增大。最终，它们作为与夸克-反夸克束缚态相联系的共振而主导总截面。

\begin{figure}[t]
  \centering
  \includegraphics[width=0.6\textwidth]{fig18.8.png}
  \caption{能量低于 3 GeV 时反应 $e^{+}e^{-}\to$ 强子总截面的实验测量，与 3 种夸克味的微扰 QCD 预言比较。数据取自 M.~Swartz 的汇编, \emph{Phys.\ Rev.\ D} \textbf{53}, 5268 (1996)。各种结果的完整参考文献见该文。}
  \label{fig:18.8}
\end{figure}

\section{深度非弹性散射的算符分析}

我们现在把算符乘积展开应用于 QCD 硬散射过程的另一个例子，即深度非弹性电子散射。在第~17 章中，我们发现 QCD 对深度非弹性散射的预言是精确的，但结构也很复杂。在第一个层次上，QCD 意味着深度非弹性散射由部分子模型描述，其中入射电子与携带质子总动量一部分的夸克和反夸克散射。这些份额由部分子分布函数确定，后者反映质子波函数的形式并由软的 QCD 动力学确定。然而，我们在第~17.5 节看到，QCD 微扰论的效应使部分子分布的形式随动量转移 $Q^2$ 而改变。我们现在将证明，这幅图景的很大部分可以从我们的新观点出发、利用算符乘积展开重构出来。

在前一节中，我们分三步推导了 $e^{+}e^{-}$ 湮灭截面的 OPE 关系。首先，我们用光学定理把这个截面与流的乘积的矩阵元联系起来。其次，我们把算符乘积展开应用于流的乘积。不幸的是，这个展开只能在非物理的运动学区域中使用。然而，在第三步中，我们用色散关系的方法把这个非物理结果与我们希望预言的截面的积分联系起来。在我们对深度非弹性散射的讨论中，我们将经历同样的三步。为了得到最终结果，我们将需要增加第四步，涉及 QCD 算符重标度。

\subsection{深度非弹性散射的运动学}

我们首先写下深度非弹性散射截面的普遍表达式。深度非弹性电子散射到末态 $f$ 的矩阵元按图~\ref{fig:18.9}(a) 所示计算：
\begin{equation}
  \mathcal{M} \propto \int d^4x\ e^{iq\cdot x}\ \langle f|J^\mu(x)|P\rangle ,
  \label{eq:18.101}
\end{equation}
其中 $J^\mu(x)$ 是夸克电磁流。这个表达式的核心是流在质子与某个高能强子态之间的强子矩阵元。这个矩阵元必须取平方并对可能的末态求和。这个求和可以用光学定理，通过把它与质子态中两个流的前向矩阵元联系起来而计算，如图~\ref{fig:18.9}(b) 所示。定义
\begin{equation}
  W^{\mu\nu}(P, q) = \int \frac{d^4x}{2\pi}\ e^{iq\cdot x}\
    \langle P|J^\mu(x)J^\nu(0)|P\rangle ,
  \label{eq:18.102}
\end{equation}
对质子的自旋取平均。这个量称为\emph{前向 Compton 振幅}，因为如果它在 $q^2 = 0$ 处求值并与物理极化矢量收缩，它给出光子-质子散射的前向振幅：
\begin{equation}
  i\mathcal{M}(\gamma p \to \gamma p) = (ie)^2\ \epsilon^{*}_\mu(q)
    \epsilon_\nu(q)\ \bigl( -iW^{\mu\nu}(P, q) \bigr) .
  \label{eq:18.103}
\end{equation}
然而，在下面的讨论中，我们将需要对一般的类空 $q$ 和一般的极化态分析 \eqref{eq:18.102}。

\begin{figure}[t]
  \centering
  \includegraphics[width=0.6\textwidth]{fig18.9.png}
  \caption{深度非弹性电子散射截面的计算：(a) 振幅的一般结构；(b) 光学定理的应用。}
  \label{fig:18.9}
\end{figure}

质子上的 Compton 散射的光学定理为
\begin{equation}
  \frac{1}{2}\sum_f \int d\Pi_f\ |\mathcal{M}(\gamma p \to f)|^2
  = 2\,\mathrm{Im}\,\mathcal{M}(\gamma p \to \gamma p) .
  \label{eq:18.104}
\end{equation}
在 \eqref{eq:7.49} 中给出的推广中，这一结果扩展到初末光子极化可以任意不同的更一般情形。把 \eqref{eq:18.104} 转写到 $W^{\mu\nu}$，我们得到
\begin{equation}
  2\,\mathrm{Im}\,W^{\mu\nu}(P, q) = \sum_f \int d\Pi_f\ \langle P|
    J^\mu(-q)|f\rangle\ \langle f|J^\nu(q)|P\rangle ,
  \label{eq:18.105}
\end{equation}
其中 $J^\mu(q)$ 是流的 Fourier 变换。

我们现在可以用 $W^{\mu\nu}$ 来计算深度非弹性截面，用 \eqref{eq:18.105} 表示最后一个因子的平方。截面应当对初始电子自旋取平均、对末态电子自旋求和。于是，
\begin{equation}
  \sigma(ep \to eX) = \int \frac{d^3k'}{(2\pi)^3}\ \frac{1}{4kk'}\ \Bigl[
    \sum_{\text{spins}} \bar{u}(k)\gamma^\mu u(k')\ \bar{u}(k')\gamma^\nu
    u(k) \Bigr]\ \bigl( \frac{e^2}{q^2} \bigr)^2\ 2\,\mathrm{Im}\,W_{\mu\nu}
    (P, q) .
  \label{eq:18.106}
\end{equation}
电子旋量乘积可以计算为
\begin{equation}
  \sum_{\text{spins}} \bar{u}(k)\gamma^\mu u(k')\ \bar{u}(k')\gamma^\nu
    u(k) = \frac{1}{2}\mathrm{tr}\bigl[ \slashed{k}\gamma^\mu
    \slashed{k}'\gamma^\nu \bigr]
  = 2\bigl( k^\mu k'^{\nu} + k^\nu k'^{\mu} - g^{\mu\nu}k\cdot k'
    \bigr) .
  \label{eq:18.107}
\end{equation}
把对末态电子动量 $k'$ 和散射角 $\theta$ 的积分转换为对我们在第~17.3 节中引入的无量纲变量 $x$ 和 $y$ 的积分是有用的。这些变量用初末电子能量 $k$ 和 $k'$ 表示为
\begin{equation}
  x = \frac{Q^2}{2P\cdot q},\qquad
  y = \frac{P\cdot q}{P\cdot k} ,
  \label{eq:18.108}
\end{equation}
其中 $Q^2 = -q^2 = 2kk'(1-\cos\theta)$。

$W^{\mu\nu}$ 的普遍形式受 Lorentz 不变性和流守恒的约束为
\begin{equation}
  W^{\mu\nu}(P, q) = \Bigl( -g^{\mu\nu} + \frac{q^\mu q^\nu}{q^2} \Bigr)\,
    W_1(x, Q^2) + \Bigl( P^\mu - \frac{P\cdot q}{q^2}q^\mu \Bigr)\,
    \Bigl( P^\nu - \frac{P\cdot q}{q^2}q^\nu \Bigr)\,
    W_2(x, Q^2) ,
  \label{eq:18.113}
\end{equation}
其中 $W_1$ 和 $W_2$ 是标量结构函数。深度非弹性散射的截面可以用这些结构函数来表达。在部分子模型中，$W_1$ 和 $W_2$ 用部分子分布函数给出。关系为
\begin{equation}
  W_2(x, Q^2) = \sum_f Q_f^2\,f_f(x, Q^2),\qquad
  2x\,W_1(x, Q^2) = W_2(x, Q^2) ,
  \label{eq:18.120}
\end{equation}
具有自由费米子特征的 $y$ 依赖，如式~\eqref{eq:17.125} 中那样。最后，从 \eqref{eq:18.120} 代入 $W_2$ 的虚部，我们精确地恢复这个部分子模型表达式：
\begin{equation}
  \frac{d^2\sigma}{dx\,dy}\bigl( ep \to eX \bigr)
  = \frac{2\pi\alpha^2}{Q^4}\ \Bigl[ 1 + (1-y)^2 \Bigr]\
    \sum_f Q_f^2\,f_f(x) .
  \label{eq:18.123}
\end{equation}
这个方程将为我们提供一个参照点，以便与我们随着分析的继续而导出的更一般表达式进行比较。

\subsection{算符乘积的展开}

由于前向 Compton 振幅是流的乘积的矩阵元，计算 $W^{\mu\nu}$ 的一个替代策略是把这一乘积展开为局域算符的级数。与部分子模型的计算一样，这个方法利用渐近自由。然而，在这种情况下，这一假定被更直接地应用。算符乘积系数的计算将显式地发生在量级为 $1/Q$ 的小距离处，因此我们可以在一个耦合常数为 $\alpha_s(Q^2)$ 的微扰论中计算这些系数。

在前一节中，我们通过考虑乘积中夸克场的各种收缩方式来计算对流的乘积的真空期望值有贡献的算符的系数。在这里，我们应该注意算符 1 对 Compton 散射振幅没有贡献。领头贡献来自能够在质子波函数中产生和湮灭夸克的算符。

两个流算符乘积中最重要的项来自同味夸克的夸克流 $\bar{q}_f\gamma^\mu q_f$ 的乘积。因此我们将从研究单个夸克流的 OPE 开始。在 $\alpha_s$ 的零阶，夸克流的算符乘积的领头项由下式给出
\begin{equation}
  \bar{q}\gamma^\mu q(x)\ \bar{q}\gamma^\nu q(0)
  = \bar{q}(x)\gamma^\mu q(x)\ \bar{q}(0)\gamma^\nu q(0)
    + \text{(contractions)} ,
  \label{eq:18.124}
\end{equation}
其中收缩应当作为夸克场的 Feynman 传播子来求值。具有显式收缩的项随着 $x\to 0$ 是奇异的；其余项是非奇异的，因此在短距离极限中不那么重要。在具有不同味夸克的流的 OPE 中，没有相应的奇异项；我们将在下面论证，这个结论即使在 $\alpha_s$ 领头阶之外也成立。

为了求 $W^{\mu\nu}$，我们必须如 \eqref{eq:18.102} 所示对 \eqref{eq:18.124} 中的各项取 Fourier 变换。当我们这样做时，应该记住传播子不仅携带 Fourier 变换动量 $q$，还携带通过夸克场传入的任何动量。为了计及这一点，方便的做法是表示
\begin{equation}
  \int d^4x\ e^{iq\cdot x}\ \bar{q}(x)\gamma^\mu q(x)\ \cdots
  = \int d^4x\ e^{iq\cdot x}\ \bar{q}(0)\,e^{ix\cdot d}\,\gamma^\mu\,
    q(0)\ \cdots ,
  \label{eq:18.125}
\end{equation}
其中导数 $d$ 向右作用在夸克场上。注意这个贡献具有图~\ref{fig:18.10} 右边第一个图形的结构。类似地，\eqref{eq:18.124} 中所示的第二个收缩具有图~\ref{fig:18.10} 中第二个图形的形式。

\begin{figure}[t]
  \centering
  \includegraphics[width=0.55\textwidth]{fig18.10.png}
  \caption{两个流的算符乘积在短距离极限下的展开。}
  \label{fig:18.10}
\end{figure}

在短距离极限下，动量 $q$ 将大于进入夸克场的任何外动量。因此我们应该展开
\begin{equation}
  \frac{1}{(d+q)^2} = \frac{1}{Q^2 - 2iq\cdot d + d^2}
  = \frac{1}{Q^2}\ \sum_{n=0}^{\infty}\ \Bigl( \frac{2iq\cdot d - d^2}
    {Q^2} \Bigr)^n .
  \label{eq:18.126}
\end{equation}
我们将在下面论证，分子中含 $d^2$ 的项不重要，可以舍弃。然而，我们应该保留比值 $(2iq\cdot d/Q^2)$ 的所有幂次。这个比值分母中有 $Q^2$，因此在短距离极限下形式上是压低的。然而，在部分子模型中
\begin{equation}
  \frac{2P\cdot q}{Q^2} = \frac{1}{x} ,
  \label{eq:18.127}
\end{equation}
所以，最终所有这些项必须同等重要。我们稍后将看到这是如何运作的。

把算符乘积 \eqref{eq:18.124} 化简为有用形式所需的最后一步是化简 Dirac 矩阵的乘积。我们从 \eqref{eq:18.113} 知道，在对质子自旋取平均后，$W^{\mu\nu}$ 将在 $\mu$ 和 $\nu$ 的交换下对称。因此，对 OPE 做对称化没有害处。于是我们可以利用恒等式把三个 Dirac 矩阵的乘积化简为一个
\begin{equation}
  \frac{1}{2}\bigl( \gamma^\mu\gamma^\alpha\gamma^\nu + \gamma^\nu
    \gamma^\alpha\gamma^\mu \bigr) = g^{\mu\alpha}\gamma^\nu + g^{\nu\alpha}
    \gamma^\mu - g^{\mu\nu}\gamma^\alpha ,
  \label{eq:18.128}
\end{equation}
这个恒等式很容易从反对易关系证明。利用 \eqref{eq:18.126} 和 \eqref{eq:18.128}，我们可以把 \eqref{eq:18.125} 改写为
\begin{equation}
  \int d^4x\ e^{iq\cdot x}\ \bar{q}(x)\gamma^\mu q(x)\ \bar{q}(0)\gamma^\nu
    q(0) = \bar{q}(0)\ \Bigl[ \gamma^\mu(i d^\nu) + \gamma^\nu(i d^\mu)
      - g^{\mu\nu}\slashed{d} \Bigr]\ \sum_{n=0}^{\infty}\ \frac{(2iq\cdot d)^n}
        {Q^{2n+2}}\ q(0) ,
  \label{eq:18.129}
\end{equation}
我们可以去掉项 $(\slashed{d})q$，它在 $\alpha_s$ 的领头阶消失，因为夸克场满足 Dirac 方程。为了计算 $W_1$ 和 $W_2$，我们也可以舍弃具有显式 $q^\mu$ 因子的项，因为它们在最终将被整理成普遍形式 \eqref{eq:18.113}。于是，最终 \eqref{eq:18.125} 取如下形式
\begin{equation}
  \int d^4x\ e^{iq\cdot x}\ \bar{q}(x)\gamma^\mu q(x)\ \bar{q}(0)\gamma^\nu
    q(0) = -\frac{1}{Q^2}\ \sum_{n=0}^{\infty}\ \Bigl( \frac{2i}{Q^2}
      \Bigr)^n\ \bar{q}(0)\ \gamma^{(\mu}q^{\mu_1}\cdots q^{\mu_n)}
      \gamma^{\nu)}\ (id_{\mu_1})\cdots(id_{\mu_n})\,q(0) ,
  \label{eq:18.130}
\end{equation}
在 $\mu\leftrightarrow\nu$ 下对称化。

\eqref{eq:18.124} 中的第二项与第一项的差别在于交换点 $x$ 和 0 以及交换指标 $\mu$ 和 $\nu$。因此它的 Fourier 变换由 \eqref{eq:18.130} 作替换 $q\to -q$ 给出。因此完整的算符乘积只包含 $q$ 的偶函数项。来自算符乘积奇异项的所有其余贡献都包含算符
\begin{equation}
  \bar{q}\ (iD^{\mu_1})\cdots(iD^{\mu_k})\,q ,
  \label{eq:18.131}
\end{equation}
其指标数为偶数，这些指标要么与 $\mu$ 或 $\nu$ 等同，要么与 $q$ 的幂次收缩。为了写出算符乘积展开的相关项，我们将以两种方式修改这个算符。首先，由于 \eqref{eq:18.131} 中的算符有 $n$ 个矢量指标，它包含在 Lorentz 群的许多不同不可约表示下变换的分量。每个分量在重整化下有不同的重标度律。然而，我们将在下面看到，只有 \eqref{eq:18.131} 中自旋最高的分量与我们的分析相关。这个分量通过对指标 $\mu_1, \ldots, \mu_n$ 完全对称化，然后减去正比于 $g_{\mu_i\mu_j}$ 的项、使算符对任意指标对无迹而得到。当我们写出流的算符乘积时，我们将只保留这个分量。其次，算符 \eqref{eq:18.131} 在规范变换下不能简单地变换。由于原始流 $\bar{q}\gamma^\mu q$ 对色规范变换是不变的，两个流的算符乘积必须是规范不变算符的和。我们可以把 \eqref{eq:18.131} 中的每个因子 $(id_{\mu})$ 替换为协变导数 $(iD_\mu)$ 来使它规范不变。这一修改只增加正比于强耦合常数 $g$ 的项，因此对我们推导算符乘积系数没有影响。

综合这些改动，让我们如下定义具有味 $f$ 夸克的自旋-$n$ 算符：
\begin{equation}
  \mathcal{O}^{f}_{\mu_1\cdots\mu_n} = \bar{q}_f\ \gamma^{(\mu_1}\
    (iD^{\mu_2})\cdots(iD^{\mu_n})\,q_f - \text{traces} ,
  \label{eq:18.132}
\end{equation}
指标对称化并有适当的减除。我们可以用这些算符写出两个流 $J^\mu$ 的 OPE 最奇异部分的最终表达式。这个算符乘积中的领头项来自 \eqref{eq:18.130} 和相应的 $q\to -q$ 收缩。提取这些表达式中包含最高自旋算符 \eqref{eq:18.132} 的部分，我们得到
\begin{equation}
  J^\mu(x)J^\nu(0) \sim \sum_f \sum_{n\ \text{even}} C_n(q^2)\,
    \mathcal{O}^{f\,\mu\nu}_{\mu_1\cdots\mu_n}\ q^{\mu_1}\cdots q^{\mu_n}
    + \cdots ,
  \label{eq:18.133}
\end{equation}
其中对 $n$ 的求和只对偶数整数进行。

表达式 \eqref{eq:18.133} 是在 $\alpha_s$ 的领头阶导出的。高阶费曼图将对系数函数贡献 $\alpha_s(Q^2)$ 阶的修正。只有当这些修正被大的对数相乘时，它们才重要。如果我们把右边出现的算符 $\mathcal{O}^{(n)}$ 视为在重整化标度 $Q$ 处归一化，就没有大的动量比可用来增强对系数函数的 QCD 修正。大的对数修正仍可能在计算的稍后阶段出现，即当我们计算算符 $\mathcal{O}^{(n)}$ 的矩阵元时。

从展开 \eqref{eq:18.133}，通过在质子态中取它的期望值，直截了当地计算出 $W^{\mu\nu}$ 的展开。为了进行这一计算，我们需要知道算符 $\mathcal{O}^{(n)}$ 的质子矩阵元。注意这些矩阵元不可能依赖于动量 $q$ 的方向，因为那种依赖已被分离到系数函数中。这意味着只有质子动量 $P$ 可用于承载矩阵元的矢量指标。因此我们可以把 $\mathcal{O}^{(n)}$ 的自旋平均矩阵元写为
\begin{equation}
  \langle P|\mathcal{O}^{(n)}_{\mu_1\cdots\mu_n}|P\rangle = A_n^f\cdot
    2\,P_{\mu_1}\cdots P_{\mu_n} - \text{traces} .
  \label{eq:18.134}
\end{equation}
系数 $A_n^f$ 是无量纲的。它们不完全是纯数，因为它们依赖于算符的重整化标度，但在接下来的几段中我们将把它们视为常数。

对 $n = 1$ 的情形，算符 $\mathcal{O}^{(1)}$ 简单地化简为夸克味流 $\bar{q}_f\gamma^\mu q_f$；在这种情况下，算符的归一化与任何标度无关，系数 $A_1^f$ 是真正的常数。从我们在第~6.2 节对形状因子的普遍讨论中，我们知道守恒味流在零动量转移处的质子矩阵元由下式给出
\begin{equation}
  \langle P|\bar{q}_f\gamma^\mu q_f|P\rangle = \bar{u}(P)\gamma^\mu u(P)\,
    F_{f1}(0) ,
  \label{eq:18.135}
\end{equation}
其中 $F_{f1}(0)$ 等于质子在态中相应的守恒荷的值。对夸克流，这个荷就是态 $|P\rangle$ 中味 $f$ 的夸克数（减去反夸克数），我们将称之为 $N_f$。对质子自旋平均 \eqref{eq:18.135}，我们得到
\begin{equation}
  \langle P|\bar{q}_f\gamma^\mu q_f|P\rangle = 2P^\mu\cdot N_f .
  \label{eq:18.136}
\end{equation}
因此，对 $n = 1$，
\begin{equation}
  A_1^f = N_f = \begin{cases}2 & u\\ 1 & d\\ 0 & \text{other}.\end{cases}
  \label{eq:18.137}
\end{equation}
类似地，$\mathcal{O}^{(2)f}$ 是夸克味 $f$ 对 QCD 能量-动量张量的贡献：
\begin{equation}
  T^{\mu\nu} = \sum_f \mathcal{O}^{(2)f\,\mu\nu} + \cdots .
  \label{eq:18.138}
\end{equation}
因此，$A_2^f$ 是质子总能量-动量中由夸克味 $f$ 携带的份额。

当我们用 \eqref{eq:18.133} 和算符矩阵元的表达式 \eqref{eq:18.134} 来求 $W^{\mu\nu}$ 的级数时，我们得到
\begin{equation}
  W^{\mu\nu} = \sum_f \Bigl[ \sum_{n\ \text{even}} C_n(q^2)\ A_n^f\ \cdots
    \Bigr] .
  \label{eq:18.139}
\end{equation}
完成求和后，人们恢复结构函数的部分子模型表达式，其中部分子分布函数由算符矩阵元的适当矩给出。

\subsection{矩求和规则}

算符乘积展开给出了计算部分子模型高阶扭转修正、以及理解结构函数重整化群演化的系统方法。关键对象是 twist-2 算符的矩阵元。这些矩阵元满足重整化群方程，结构函数矩的重整化群演化由 twist-2 算符的反常维数给出。

矩求和规则把部分子分布函数的矩与算符矩阵元 $A_n^f$ 联系起来。使用色散关系，人们导出
\begin{equation}
  \int_0^1 dx\ x^{n-1}\,f_f^+(x, Q^2) = A_n^f ,
  \label{eq:18.154}
\end{equation}
对偶数的 $n$。这些关系称为深度非弹性形状因子的\emph{矩求和规则}。它们把由式~(18.146) 定义的部分子分布函数的 $x$ 矩与 twist-2 算符的质子矩阵元联系起来。

因为 $W_2$ 是 $\nu$ 的对称函数，矩求和规则只对偶数的 $n$ 成立。然而，在深度非弹性中微子散射中，$W^{\mu\nu}$ 中有第三个形状因子，与弱相互作用流的矢量部分和轴矢量部分之间的干涉项相联系。在习题~18.2 中，我们证明这个形状因子可以用来导出奇数 $n$ 的一组求和规则：
\begin{equation}
  \int_0^1 dx\ x^{n-1}\,f_f^-(x, Q^2) = A_n^f ,
  \label{eq:18.155}
\end{equation}
其中 $A_n^f$ 是奇数 $n$ 时质子矩阵元 \eqref{eq:18.134} 的系数，$f_f^-(x)$ 是一个形状因子，在部分子模型中它求值为
\begin{equation}
  f_f^-(x) = f_f(x) - f_{\bar{f}}(x) .
  \label{eq:18.156}
\end{equation}
把这一信息与 \eqref{eq:18.136} 下方给出的论证结合起来，我们可以看到从深度非弹性形状因子定义部分子分布函数具有正确的归一化。利用 \eqref{eq:18.137}，我们得到
\begin{equation}
  \int_0^1 dx\ f_f^-(x) = N_f ,
  \label{eq:18.157}
\end{equation}
即质子中味 $f$ 夸克的（净）数目。类似地，\eqref{eq:18.154} 和 \eqref{eq:18.138} 蕴含
\begin{equation}
  \int_0^1 dx\ x\,f_f^+(x) = \langle x\rangle_f ,
  \label{eq:18.158}
\end{equation}
其中 $\langle x\rangle_f$ 是质子总能量-动量中由味 $f$ 的夸克和反夸克携带的份额。

\subsection{算符重标度}

如果系数 $A_n^f$ 是真正的常数，关系 \eqref{eq:18.154} 和 \eqref{eq:18.155} 将与满足 Bjorken 标度的部分子分布函数一致。然而，正如我们在 \eqref{eq:18.134} 下方所评述的，这些因子实际上依赖于 $Q^2$，因为这是算符乘积展开 \eqref{eq:18.133} 中算符的归一化点。由于这种依赖只通过算符重标度而来，它只涉及 $Q^2$ 的对数，因此只贡献对 Bjorken 标度的缓慢破坏。通过对 twist-2 算符矩阵元的领头对数修正求和，我们可以定量地算出部分子分布函数对 $Q^2$ 的依赖。

为了计及这些修正，让我们首先（正如我们下面将看到的，这是不正确的）假定 twist-2 算符 \eqref{eq:18.132} 在没有算符混合的情况下重整化。于是算符 $\mathcal{O}^{(n)}$ 矩阵元的领头对数修正可以通过把在 $Q$ 处归一化的算符重标度到在量级为 1 GeV 的标准参考点 $\mu$ 处归一化的算符来求和。这些约定之间的关系为
\begin{equation}
  \bigl[ \mathcal{O}^{(n)} \bigr]_Q = \Bigl( \frac{\alpha_s(Q^2)}
    {\alpha_s(\mu^2)} \Bigr)^{a_n^{qq}/2b_0}\ \bigl[ \mathcal{O}^{(n)}
      \bigr]_\mu ,
  \label{eq:18.159}
\end{equation}
其中 $a_n^{qq}$ 是 $\mathcal{O}^{(n)}$ 的 $\gamma$ 函数的第一个系数。于是因子 $A_n^f$ 将按如下方式依赖于 $Q^2$
\begin{equation}
  A_n^f(Q^2) = \Bigl( \frac{\alpha_s(Q^2)}{\alpha_s(\mu^2)} \Bigr)^{
    a_n^{qq}/2b_0}\ A_n^f(\mu^2) .
  \label{eq:18.160}
\end{equation}
这个方程与 \eqref{eq:18.76} 中写出的算符乘积系数的标度依赖一致，适用于流的算符乘积的特殊情形 $a_1 = a_2 = 0$。为了找到重标度因子的显式形式，我们必须计算 $a_n^{qq}$。

\begin{figure}[t]
  \centering
  \includegraphics[width=0.5\textwidth]{fig18.13.png}
  \caption{贡献夸克 twist-2 算符反常维数的图形。}
  \label{fig:18.13}
\end{figure}

为了计算夸克 twist-2 算符的 $\gamma$ 函数，我们必须计算它们的算符重标度抵消项。它们由图~\ref{fig:18.13} 中所示图形确定。在外部动量 $p$ 通过夸克线进入、且注入算符的外动量为零的条件下计算这些图形就足够了。在这些条件下，算符 $\mathcal{O}^{(n)}$ 的矩阵元在领头阶等于
\begin{equation}
  \langle p|\mathcal{O}^{(n)}_{\mu_1\cdots\mu_n}|p\rangle
  = \bar{u}(p)\,\gamma_{\mu_1}\,p_{\mu_2}\cdots p_{\mu_n}\,u(p) .
  \label{eq:18.161}
\end{equation}
在此处以及在讨论中所有稍后的位置，我们将把 $\mathcal{O}^{(n)}$ 的矩阵元视为在 $n$ 个指标上对称化并减除了所有可能的迹。我们现在必须计算图~\ref{fig:18.13} 中的图形，并收集所有重标度这个结构的项。

图~\ref{fig:18.13} 中的第一个图形计算起来相当直截了当。我们利用恒等式 \eqref{eq:6.40} 合并分母：
\begin{equation}
  \frac{1}{k^2(k-p)^2} = \int_0^1 dx\ \frac{1}{(k-xp)^4} ,
  \label{eq:18.163}
\end{equation}
右边分母中的量是 $\bar{k} = k - xp$ 和 $\Delta = -x(1-x)p^2$。我们现在必须移动积分，在分子中代入 $\bar{k} = k + xp$，并挑出正比于 $p$ 的 $(n-1)$ 次幂的项。如果这一项含有因子 $g^{\mu_i\mu_j}$，我们可以舍弃它，因为它对更高扭转算符的系数有贡献，并且无论如何，当我们减除迹时它会被去掉。因此，当我们把其他因子换成 $(xp)$ 时，我们必须仔细选择两个 $\bar{k}$ 因子用什么替换。下面这些选择，利用 $\bar{k}$ 积分的转动对称性化简，不能给出有用的贡献：
\begin{equation}
  \int \frac{d^4\bar{k}}{(2\pi)^4}\ \frac{\bar{k}^\mu\bar{k}^\nu}
    {(\bar{k}^2-\Delta)^4} = \frac{i}{(4\pi)^2}\ \frac{g^{\mu\nu}}{4}\
      \int d\bar{k}\ \frac{\bar{k}^2}{(\bar{k}^2-\Delta)^4} .
  \label{eq:18.164}
\end{equation}
剩下的一个 $\bar{k}$ 因子放置方式是
\begin{equation}
  \bar{k}^{\mu_1}\bar{k}^{\mu_n}\gamma^{\mu_1}\gamma^{\mu_n}
  \to \frac{1}{4}\bar{k}^2\,\gamma^{\mu_1}\gamma^{\mu_2}\ \cdots .
  \label{eq:18.165}
\end{equation}

因此式~(18.162) 的值为
\begin{equation}
  \text{(first diagram)} = -\frac{2}{3}\,n(n+1)\ \frac{g^2}{(4\pi)^2}\
    \frac{1}{n(n+1)}\,\Gamma(2-d/2)\ \bar{u}(p)\,\gamma^{\mu_1}\,p^{\mu_2}
    \cdots p^{\mu_n}\,u(p)\ \cdots .
  \label{eq:18.166}
\end{equation}
对算符 $\mathcal{O}^{(n)}$ 的重标度还有额外的贡献，这一点并不那么明显。然而，注意 \eqref{eq:18.132} 中的协变导数包含规范场的显式因子，
\begin{equation}
  iD_\mu = i\partial_\mu - gA^a_\mu t^a ,
  \label{eq:18.167}
\end{equation}
这些因子可以与外腿上的规范场顶角收缩。这些贡献产生图~\ref{fig:18.13} 中的第二和第三个图形。\eqref{eq:18.167} 中两个 $A^a_\mu$ 因子相互收缩的项正比于 $G^{\mu\nu}$，因此对领头扭转算符的重标度没有贡献。

我们刚才描述的贡献具有对 $j$ 求和的形式，其中 $\mu_j$ 是包含收缩的导数的指标。于是图~\ref{fig:18.13} 中第二个图形是下列积分对 $j$ 的和：
\begin{equation}
  \text{(second diagram)} = \sum_j\ \int \frac{d^4k}{(2\pi)^4}\ \frac{1}{k^2}
    \ \frac{1}{(k-p)^2}\ \bar{u}(p)\,\gamma^\mu\ \slashed{k}\ \gamma^{\mu_1}
    \cdots \cdots \,u(p) .
  \label{eq:18.168}
\end{equation}
由于 $\mu_j$ 和 $\mu_1$ 被对称化，我们可以用 \eqref{eq:18.128} 把 $\gamma^\mu\slashed{k}\gamma^{\mu_1} + \gamma^{\mu_1}\slashed{k}\gamma^\mu$ 改写为方便的形式。完成积分后，人们得到反常维数系数 $a_n^{qq}$。把这些与图~\ref{fig:18.13} 中第三个图形的贡献结合，人们得到 twist-2 算符的完整单圈反常维数矩阵。

这一计算的结果是反常维数矩阵
\begin{equation}
  \gamma^{(n)} = \frac{g^2}{8\pi^2}\ \begin{pmatrix}
    a_n^{qq} & a_n^{qg}\\ a_n^{gq} & a_n^{gg}
  \end{pmatrix} ,
  \label{eq:18.181}
\end{equation}
其中系数由下式给出
\begin{equation}
  \begin{aligned}
    a_n^{qq} &= \frac{4}{3}\Bigl[ 1 - \frac{2}{n(n+1)} + 4\sum_{j=2}^{n}
      \frac{1}{j} \Bigr],\\
    a_n^{qg} &= \frac{n^2+n+2}{n(n+1)},\\
    a_n^{gq} &= \frac{4}{3}\ \frac{n^2+n+2}{n(n^2-1)},\\
    a_n^{gg} &= 6\Bigl[ 1 - \frac{2}{n(n-1)} - \frac{2}{(n+1)(n+2)}
      + 4\sum_{j=2}^{n}\frac{1}{j} \Bigr] - \frac{2}{3}n_f .
  \end{aligned}
  \label{eq:18.181b}
\end{equation}
注意这个矩阵不对称。在最后一行中，$n_f$ 是夸克味的数目，在此情形下等于 1；这一项来自 \eqref{eq:16.74}。

在现实的情形中，QCD 包含几种夸克味——$u$、$d$、$s$，以及当我们在足够大的动量下工作、可以忽略这些粒子质量时的 $c$ 和 $b$。于是反常维数矩阵 $\gamma^{(n)}$ 的大小为 $(n_f+1)\times(n_f+1)$。作用在夸克算符上的子矩阵是对角的，所有对角元由 \eqref{eq:18.181} 中的 $a_n^{qq}$ 给出。夸克-胶子和胶子-夸克元分别由 $a_n^{qg}$ 和 $a_n^{gq}$ 给出，并且与味无关。胶子对角元由 \eqref{eq:18.181} 中取现实 $n_f$ 值的 $a_n^{gg}$ 给出。这意味着胶子算符只与夸克算符的一个线性组合混合：
\begin{equation}
  \sum_f \mathcal{O}^{(n)}_f ,
  \label{eq:18.182}
\end{equation}
正交的线性组合被简单地重标度，其指数由 $a_n^{qq}$ 或式~(18.172) 给出。

让我们现在把这一算符混合分析应用于矩求和规则的求值。对奇数的 $n$，没有算符混合，因此 \eqref{eq:18.155} 右边对 $Q^2$ 的依赖由简单的重标度 \eqref{eq:18.160} 正确地给出。

对偶数的 $n$，我们必须把算符混合考虑进去。求和规则 \eqref{eq:18.154} 的右边是在标度 $Q$ 处归一化的 twist-2 算符的质子矩阵元。让我们把这些算符的任意线性组合写为
\begin{equation}
  c^i\ \bigl[ \mathcal{O}^{(n)}_i \bigr]_Q ,
  \label{eq:18.183}
\end{equation}
其中指标 $i$ 遍历 $g$ 和各种味 $f$。为了把这个算符重标度到固定的参考动量 $\mu$，我们把系数改写为 $\gamma^{(n)}$ 的左本征矢基，并按照 \eqref{eq:18.159} 重标度每个本征矢。用重标度系数矩阵 $a^n$ 表示，我们可以把重标度抽象地写为
\begin{equation}
  c^i\ \bigl[ \mathcal{O}^{(n)}_i \bigr]_Q \to c^i\ \Bigl( \frac{\alpha_s(Q^2)}
    {\alpha_s(\mu^2)} \Bigr)^{a^n/2b_0}\ \bigl[ \mathcal{O}^{(n)}_i
      \bigr]_\mu .
  \label{eq:18.184}
\end{equation}
这个重标度，其中 $c^i$ 作用在矩阵 $(a^n)$ 的左边，正是式~\eqref{eq:18.79} 所要求的方案。

让我们对 $n = 2$ 的情形显式地算出这一点。矩求和规则 \eqref{eq:18.154} 的右边由 $\mathcal{O}^{(2)}_f$ 的矩阵元给出。我们把它改写为
\begin{equation}
  \sum_f \mathcal{O}^{(2)}_f = \Bigl( \sum_f \mathcal{O}^{(2)}_f \Bigr)\
    + \cdots ,
  \label{eq:18.185}
\end{equation}
其中第一个组合是能量-动量张量，第二项与胶子算符 $\mathcal{O}^{(2)}_g$ 混合。对 $n = 2$ 作用在 $(\mathcal{O}^{(2)}_f, \mathcal{O}^{(2)}_g)$ 上的反常维数矩阵具有系数
\begin{equation}
  \begin{pmatrix}
    a_2^{qq} & n_f\,a_2^{qg}\\
    a_2^{gq} & a_2^{gg}
  \end{pmatrix}
  = \begin{pmatrix}
    \frac{16}{9} & \frac{2}{3}n_f\\
    \frac{16}{9} & -\frac{2}{3}n_f
  \end{pmatrix} .
  \label{eq:18.186}
\end{equation}
这个矩阵的左本征矢及其相应的本征值为
\begin{equation}
  (1, 1):\ a = 0;\qquad
  \Bigl( \frac{4}{3}, -\frac{n_f}{2} \Bigr):\ a = \frac{4}{3}\Bigl(
    \frac{16}{9} + \frac{2}{3}n_f \Bigr) .
  \label{eq:18.187}
\end{equation}
注意第一个本征矢给出反常维数为零的算符组合 $\sum_f\mathcal{O}^{(2)}_f$。这个算符实际上就是 QCD 的总能量-动量张量，
\begin{equation}
  T^{\mu\nu} = \sum_f \mathcal{O}^{(2)f\,\mu\nu} ,
  \label{eq:18.188}
\end{equation}
它必有 $\gamma = 0$。如果我们把 \eqref{eq:18.185} 中的第二项按分量 \eqref{eq:18.187} 展开，我们就可以计算算符重标度的完整形式。我们得到
\begin{equation}
  \bigl[ \mathcal{O}^{(2)}_f \bigr]_Q = \frac{1}{n_f}\ \Bigl[
    T^{\mu\nu} - \frac{4}{3}\frac{n_f}{16/3+n_f}\ \Bigl( \frac{\log(Q^2/
      \Lambda^2)}{\log(\mu^2/\Lambda^2)} \Bigr)^{-\lambda_2/2b_0}\ \cdots
    \Bigr] ,
  \label{eq:18.189}
\end{equation}
其中 $T^{\mu\nu}$ 是能量-动量张量 \eqref{eq:18.188}。$n = 2$ 矩求和规则的右边由这个算符的质子矩阵元的系数给出。为了计算这个系数，我们需要定义 $A_n^f$ 的胶子类比，与 \eqref{eq:18.134} 类似地写出
\begin{equation}
  \langle P|\mathcal{O}^{(n)g}_{\mu_1\cdots\mu_n}|P\rangle
  = A_n^g\cdot 2\,P_{\mu_1}\cdots P_{\mu_n} - \text{traces} .
  \label{eq:18.190}
\end{equation}
对 $n = 2$ 的情形，我们特别注意到
\begin{equation}
  \langle P|T^{\mu\nu}|P\rangle = 2P^\mu P^\nu ,
  \label{eq:18.191}
\end{equation}
因此，\eqref{eq:18.188} 蕴含
\begin{equation}
  \sum_f A_2^f + A_2^g = 1 .
  \label{eq:18.192}
\end{equation}
如果我们把 \eqref{eq:18.189} 中的每个算符替换为相应的系数 $A_2^f(\mu)$，我们将得到 $n = 2$ 矩求和规则右边的一个表达式，它使其对 $Q^2$ 的依赖显式化。

尽管表达式 \eqref{eq:18.189} 相当复杂，它在极端极限 $Q^2\to\infty$ 下具有简单的形式。在渐近 $Q^2$ 处，\eqref{eq:18.189} 的最后两项趋于零，\eqref{eq:18.189} 的右边变为一个固定数乘以能量-动量张量。于是，利用 \eqref{eq:18.191}，我们可以完全求出 $n = 2$ 矩求和规则：
\begin{equation}
  \int_0^1 dx\ x\,f_f^+(x, \infty) = \frac{1}{n_f}\ .
  \label{eq:18.193}
\end{equation}
在这个极端极限下，我们发现每种夸克味携带质子能量-动量的同一个固定份额。由 \eqref{eq:18.192}，其余部分由胶子携带。举例说明，在 $n_f = 4$ 的理论中，每种夸克味携带质子总动量的 $3/28$，胶子携带剩余的 $4/7$。图~\ref{fig:18.16} 显示了从现实的部分子分布出发，这些渐近结果被逼近得多么缓慢。

\begin{figure}[t]
  \centering
  \includegraphics[width=0.6\textwidth]{fig18.16.png}
  \caption{质子的各种部分子种类所携带的总能量-动量份额，作为 $Q$ 的函数，按图~\ref{fig:17.6} 所述对深度非弹性散射数据的 CTEQ 拟合。曲线的 $Q$ 依赖由 QCD 演化方程计算。}
  \label{fig:18.16}
\end{figure}

\subsection{与 Altarelli-Parisi 方程的关系}

刚才描述的算符混合分析给出对部分子分布矩的预言，这些预言意味着这些积分依赖于 $Q^2$。在不涉及算符混合的各种矩积分中，只有给出质子味量子数的 $n = 1$ 积分作为 $Q^2$ 的函数是常数。其余的作为 $\log Q^2$ 的幂次减小。类似地，$n = 2$ twist-2 算符矩阵元的一个线性组合随 $Q^2$ 保持为常数。这一关系由求和规则 \eqref{eq:18.192} 表达。为了更清楚地写出这一关系，让我们引入胶子的部分子分布，把它作为一个满足下列关系的光滑函数
\begin{equation}
  \int_0^1 dx\ x\,f_g(x, Q^2) = A_2^g(Q^2) .
  \label{eq:18.194}
\end{equation}
于是 \eqref{eq:18.192} 就变成部分子分布的总动量求和规则 \eqref{eq:17.39}：
\begin{equation}
  \int_0^1 dx\ x\ \Bigl[ \sum_f f_f^+(x, Q^2) + f_g(x, Q^2) \Bigr] = 1 .
  \label{eq:18.195}
\end{equation}
不难验证，对 $n > 2$，反常维数系数矩阵 $a^n$ 的所有本征值都是负的。因此，所有更高的矩求和规则都减小，但要服从味荷和动量守恒定律。换言之，算符重整化分析预言，随着 $\log Q^2$ 增大，部分子分布向更小的 $x$ 值移动。令人欣慰的是，这正是我们在第~17.5 节中得到的结论，我们在那里推导了 Altarelli-Parisi 方程来描述部分子分布的演化。

\problems

\textbf{18.1} 在这个习题中，我们将按照第~18.1 节的分析研究 QCD 中夸克质量参数的重整化。显式计算单圈反常维数 $\gamma_{\bar{q}q}$，验证 \eqref{eq:18.13}，并检验结果与规范参数无关。

\textbf{18.2} \emph{奇数 $n$ 的求和规则。} 在正文中，我们推导了适用于偶数 $n$ 的深度非弹性结构函数矩求和规则。在这个习题中，证明深度非弹性中微子散射中弱相互作用流的矢量部分与轴矢量部分之间的干涉项给出奇数 $n$ 的一组求和规则，并验证式~\eqref{eq:18.155}。

\textbf{18.3} \emph{胶子 twist-2 算符的反常维数。} 完成 twist-2 算符单圈反常维数矩阵的计算，验证 \eqref{eq:18.181} 中的系数 $a_n^{qg}$、$a_n^{gq}$ 和 $a_n^{gg}$。

\textbf{18.4} 在第~18.2 节对弱相互作用算符乘积展开的分析中，我们发现 QCD 修正混合两个四费米子算符。在这个习题中，求解这些算符的重整化群方程，并计算 $\Delta I = 1/2$ 和 $\Delta I = 3/2$ 振幅的有效增强因子。
